% !TEX root = % !TEX root = /Users/nai1ka/Projects/PerfX/Thesis/thesis.tex
\begin{abstract}
% skip one line to make the abstract start with indent


Each new release of a mobile application can introduce a performance regression, for example through added features or updated dependencies. On Android these regressions are difficult to detect, because user devices vary widely in hardware and operating system, so the same release can run differently from one device to another. Regression detection is well established on the server side, but Android tools usually rely on fixed thresholds that a developer sets by hand and do not account for this device variety. This thesis presents a complete, open-source system that monitors Android applications and automatically detects performance regressions between releases. Its distinguishing element is that, instead of comparing all devices together, the system groups devices with similar capabilities into performance cohorts and compares each new release against the previous one within each cohort, so that a regression appearing only on weaker devices is not hidden among the stronger ones. The system has three parts: a mobile SDK on the device, a backend, and a frontend. Moreover, it sends an alert to developers when a regression is found. Because the SDK runs on the user's device, its overhead must remain low. Under a realistic workload it added 19.5 percentage points of CPU usage, 6.3 MB (3.9\%) of memory, and 276 ms (15.9\%) of cold-start time, around four times less CPU and nine times less memory than two widely used industry SDKs. The whole system was tested end to end, which it successfully passed. These results show that post-release monitoring and automatic cohort-aware regression detection can be applied to Android applications on real devices at a low cost to the user and with little setup for the developer.

\end{abstract}